% Generated draft: requires mathematical and source review.
\documentclass[11pt]{article}
\usepackage[margin=25mm]{geometry}
\usepackage{fontspec}
\setmainfont{Latin Modern Roman}
\usepackage{amsmath,amssymb,mathtools,unicode-math}
\setmathfont{Latin Modern Math}
\usepackage{enumitem,needspace,xurl,hyperref}
\setlength{\parindent}{0pt}
\setlength{\parskip}{5pt}
\setlength{\emergencystretch}{3em}
\newcommand{\N}{\mathbb{N}}
\newcommand{\Z}{\mathbb{Z}}
\newcommand{\Q}{\mathbb{Q}}
\newcommand{\R}{\mathbb{R}}
\newcommand{\C}{\mathbb{C}}
\newcommand{\F}{\mathbb{F}}
\newcommand{\A}{\mathbb{A}}
\newcommand{\PP}{\mathbb{P}}
\newcommand{\E}{\mathbb{E}}
\newcommand{\eps}{\varepsilon}
\newcommand{\dd}{\,\mathrm{d}}
\begin{document}
\section*{UnsolvedMath ID: 30002637}
\textbf{Source problem code:} catalogue entry
\subsection{Definitions and mathematical statement}
A Ricci soliton is a Riemannian manifold \((M,g)\) equipped with a vector field \(X\) and a constant \(\lambda\in\R\) satisfying
\[
\operatorname{Ric}_g+\tfrac12\mathcal L_Xg=\lambda g.
\]
It is a gradient soliton when \(X=\nabla^g f_0\), so that \(\operatorname{Ric}_g+\nabla_g^2f_0=\lambda g\), and it is shrinking when \(\lambda>0\). An Einstein metric satisfies \(\operatorname{Ric}_g=cg\) for a constant \(c\). The supplied record does not explicitly define ``nontrivial'' or specify the manifold and operator conventions underlying its entropy statement.

Here is a precise conditional formulation in the usual closed, connected, gradient shrinking setting. Let \(M\) be a smooth, closed, connected manifold of dimension \(n\), and let \(g\) be a smooth Riemannian metric with
\[
\operatorname{Ric}_g+\nabla_g^2f_0=\frac{1}{2\tau_0}g,
\qquad \tau_0>0.
\]
For this formulation, call the soliton nontrivial when \(g\) is not Einstein. These conventions require confirmation against the original source; they are not additional hypotheses explicitly supplied by the record.

For \(\tau>0\) and \(f\in C^\infty(M)\) satisfying
\[
\int_M(4\pi\tau)^{-n/2}e^{-f}\,\dd V_g=1,
\]
define Perelman's entropy functionals by
\[
\mathcal W(g,f,\tau)=\int_M\bigl[\tau(R_g+|\nabla^g f|_g^2)+f-n\bigr](4\pi\tau)^{-n/2}e^{-f}\,\dd V_g,
\]
\[
\mu(g,\tau)=\inf_{f\text{ satisfying the normalization}}\mathcal W(g,f,\tau),
\qquad
\nu(g)=\inf_{\tau>0}\mu(g,\tau).
\]
Here \(R_g\) is scalar curvature and \(\dd V_g\) is Riemannian volume. At a soliton where the second variation is defined, write
\[
Q_g(h,k)=\left.\frac{\partial^2}{\partial s\,\partial t}\right|_{s=t=0}\nu(g+sh+tk),
\qquad h,k\in C^\infty(M,\operatorname{Sym}^2T^*M).
\]
Normalize \(f_0\) so that \(\dd\rho_0=(4\pi\tau_0)^{-n/2}e^{-f_0}\dd V_g\) has total mass one, and represent the Hessian by \(H_g\) through
\[
Q_g(h,k)=\int_M\langle H_gh,k\rangle_g\,\dd\rho_0.
\]
\textbf{Question.} For every nontrivial compact Ricci soliton in the intended source setting, does its entropy Hessian have a positive eigenvalue? Under the conventions above, this asks whether every such \((M,g,f_0,\tau_0)\) admits \(\eta>0\) and a nonzero smooth symmetric two-tensor \(h\) such that
\[
H_gh=\eta h?
\]
The positive sign is the instability convention in the supplied statement. The source record does not identify the precise operator domain, gauge restriction, or treatment of the null directions arising from diffeomorphisms and scaling. Thus the displayed realization makes the intended question explicit but remains subject to source and mathematical review.
\subsection{Short English statement}
Must every nontrivial compact Ricci soliton have a positive eigenvalue in the second variation of Perelman's \(\nu\)-entropy? The supplied record labels this question open. The precise meaning of nontriviality, the manifold conventions, and the Hessian realization require confirmation; the displayed formulation uses a closed, connected, gradient shrinking soliton with a non-Einstein metric.
\subsection{Sources}
\begin{itemize}
\item Supplied local UnsolvedMath registry record, permanent ID \(30002637\), titled ``Instability of Nontrivial Compact Ricci Solitons,'' in the Oberwolfach Reports Open Problems collection. Exact source problem code: \texttt{catalogue entry}. The proposed year is recorded as 2014; no individual proposer is supplied. The registry labels the problem \texttt{open} and the statement \texttt{corrected\_verified}; these are supplied metadata, not independent verification.
\item Supplied report citation: \emph{Analysis, Geometry and Topology of Positive Scalar Curvature Metrics} (2015), Oberwolfach Reports, pp.\ 1991--2046, DOI \texttt{10.4171/owr/2014/36}. \url{https://doi.org/10.4171/owr/2014/36}. This is also the literature-source link labelled ``original OWR report.'' No exact problem page or named speaker is supplied.
\item Supplied literature source labelled ``Kröncke stability paper'': \url{https://doi.org/10.1007/s00526-014-0748-3}. The record supplies no paper title or theorem reference.
\item Literature assessment dated exactly \texttt{2026-08-21}: supplied dated curation, not independently verified. It reports that all known nontrivial compact Ricci solitons were unstable and labels the universal instability question open. Its concluding wording refers to ``every nontrivial Ricci soliton'' without repeating compactness; the original and normalized problem statements explicitly concern compact solitons. The background instructs that time-sensitive status and cited claims require independent verification.
\end{itemize}
\end{document}
